
%%% Преамбула %%%
\usepackage[utf8]{inputenc}
\usepackage{amssymb,amsfonts,amsmath,cite,enumerate,float,indentfirst} %пакеты расширений
\usepackage[dvips]{graphicx} %вставка графики
\usepackage[main=russian,english]{babel}
\usepackage{fontspec}
\setmainfont[
BoldFont={timesnewromanbold.ttf},
ItalicFont={timesnewromanitalic.ttf},
BoldItalicFont={timesnewromanbolditalic.ttf}
]{timesnewroman.ttf}
\usepackage{bm}
\usepackage{cmap}
\usepackage{float}
\usepackage{wrapfig}
\usepackage{amsmath}
\usepackage{mathtools,cuted}
\usepackage{multicol}
\usepackage{tabstackengine}
\usepackage{fancybox,framed}
\usepackage{xpatch, nccmath}
\usepackage{lastpage}
\usepackage{totcount}
\usepackage{float}
\xpatchcmd{\NCC@ignorepar}{%
	\abovedisplayskip\abovedisplayshortskip}
{%
	\abovedisplayskip\abovedisplayshortskip%
	\belowdisplayskip\belowdisplayshortskip}
{}{}
\usepackage[onehalfspacing]{setspace} % полуторный интервал по всему документу (см. setspace)
\AtBeginDocument{%
  % Единый вертикальный интервал до и после всех display-формул (equation, align и т.д.)
  \setlength{\abovedisplayskip}{\baselineskip}%
  \setlength{\belowdisplayskip}{\baselineskip}%
  \setlength{\abovedisplayshortskip}{\baselineskip}%
  \setlength{\belowdisplayshortskip}{\baselineskip}%
  % Единый вертикальный интервал вокруг рисунков
  \setlength{\intextsep}{\baselineskip}%
  \setlength{\textfloatsep}{\baselineskip}%
  \setlength{\floatsep}{\baselineskip}%
}
\AtBeginEnvironment{figure}{\addvspace{\baselineskip}}
\stackMath
\usepackage[most]{tcolorbox}
\usepackage[figurename=Рисунок]{caption}
\usepackage{systeme}
\DeclareCaptionLabelSeparator{dblhyphen}{ -- }
\captionsetup[table]{%
  justification=raggedright,
  singlelinecheck=false,
  position=top,
  font={stretch=1},
  labelsep=dblhyphen,
  skip=0pt,
  aboveskip=0pt,
  belowskip=\baselineskip
}
\usepackage[hidelinks]{hyperref}
% \url в списке литературы и note: без «курсивного» моноширинного начертания — как основной текст
\AtBeginDocument{\urlstyle{rm}}
% Пакет для изображений
\usepackage{graphicx}
\usepackage{pdfpages} % первая страница из готового PDF (титульник)
\graphicspath{ {./images/} }
\newenvironment{Figure}
{\par\medskip\noindent\minipage{\linewidth}}
{\endminipage\par\medskip}
\makeatletter
\renewcommand{\@biblabel}[1]{#1.} % Заменяем библиографию с квадратных скобок на точку:
\makeatother
\sloppy             % Избавляемся от переполнений
\hyphenpenalty=1000000 % Частота переносов
\clubpenalty=1000000  % Запрещаем разрыв страницы после первой строки абзаца
\widowpenalty=1000000 % Запрещаем разрыв страницы после последней строки абзаца
\usepackage{changepage}% http://ctan.org/pkg/changepage
\usepackage[nopar]{lipsum}
% Отступы у страниц
\usepackage[margin=3cm]{geometry}
\geometry{left=3cm}
\geometry{right=1.5cm}
\geometry{top=2cm}
\geometry{bottom=2cm}
\parindent 1.25cm % Абзацный отступ
% Оформление библиографии и подрисуночных записей через точку
\makeatletter
\renewcommand*{\@biblabel}[1]{\hfill#1.}
\renewcommand*\l@section{\@dottedtocline{1}{1em}{1em}}
\renewcommand{\thefigure}{\arabic{figure}} % Формат рисунка секция.номер
\renewcommand{\thetable}{\arabic{table}} % Формат таблицы секция.номер
\def\redeflsection{\def\l@section{\@dottedtocline{1}{0em}{10em}}}
\usepackage{listings}
\lstset{
	language=Python,
	basicstyle=\small\ttfamily,
	columns=fullflexible,
	breaklines=true,
	breakatwhitespace=true,
	keepspaces=true,
	tabsize=4,
	showstringspaces=false
}
\usepackage{enumitem}
\setlist[itemize]{
	label=-,
	labelindent=0pt,
	leftmargin=0pt,
	labelwidth=0pt,
	labelsep=0.5em,
	itemindent=1.5cm,
	listparindent=0pt,
	topsep=0pt,
	partopsep=0pt,
	parsep=0pt,
	itemsep=0pt,
	align=left
}
\setlist[enumerate]{
	label=\arabic*),
	labelindent=0pt,
	leftmargin=0pt,
	labelwidth=0pt,
	labelsep=0.5em,
	itemindent=1.5cm,
	listparindent=0pt,
	topsep=0pt,
	partopsep=0pt,
	parsep=0pt,
	itemsep=0pt,
	align=left
}
% TOC
\addto\captionsrussian{% Replace "english" with the language you use
	\renewcommand{\contentsname}%
	{\centering  СОДЕРЖАНИЕ}%
}
\usepackage{indentfirst}
\usepackage{titlesec}
\usepackage{lipsum}
\newlength{\normalparindent}
\AtBeginDocument{\setlength{\normalparindent}{\parindent}}
\titleformat{name=\section}[block]
{\normalsize\bfseries}
{\hspace*{\normalparindent}\thesection}
{1ex}
{}
\titleformat{name=\section,numberless}[block]
{\normalsize\bfseries\centering}
{}
{0pt}
{}
\titleformat{name=\subsection}[block]
{\normalsize\bfseries}
{\hspace*{\normalparindent}\thesubsection}
{1ex}
{}
\titleformat{name=\subsection,numberless}[block]
{\normalsize\bfseries}
{}
{0pt}
{\hspace*{\normalparindent}}
\titleformat{name=\subsubsection}[block]
{\normalsize\bfseries}
{\hspace*{\normalparindent}\thesubsubsection}
{1ex}
{}
\titleformat{name=\subsubsection,numberless}[block]
{\normalsize\bfseries}
{}
{0pt}
{\hspace*{\normalparindent}}
\renewcommand\paragraph{\@startsection{paragraph}{4}{\z@}%
	% display heading, like subsubsection
	{-3.25ex\@plus -1ex \@minus -.2ex}%
	{1.5ex \@plus .2ex}%
	{\normalfont\normalsize\bfseries}}
\usepackage{tabularx}
\usepackage{booktabs}
% Формат подрисуночных записей
\usepackage{chngcntr}
\usepackage{ragged2e}
%\usepackage{times}
\newcolumntype{b}{X}
\newcolumntype{s}{>{\raggedleft\hsize=.25\hsize}X}
\newcolumntype{h}{>{\raggedleft\hsize=.5\hsize}X}
\newcolumntype{w}{>{\raggedleft\hsize=.01\hsize}X}
\newcolumntype{t}{>{\raggedleft\hsize=.03\hsize}X}
\newcolumntype{i}{>{\raggedleft\hsize=.05\hsize}X}
\newcolumntype{o}{>{\raggedleft\hsize=.14\hsize}X}
\captionsetup[figure]{%
  labelsep=dblhyphen,
  justification=centering,
  singlelinecheck=false,
  font={stretch=0.86}
}
\newcommand{\mycaption}[1]
{
	\centering
	 \caption[#1]{#1}
}
\newcommand{\imgh}[2] {
	\begin{figure}[h!]
		\center{
			\includegraphics[width=0.75\linewidth]{#1}}
		\mycaption{#2}
		\label{ris:#1}
	\end{figure}
}
\newcommand{\imghsize}[3] {
	\begin{figure}[h!]
		\center{
			\includegraphics[width={#3}\linewidth]{#1}}
		\mycaption{#2}
		\label{ris:#1}
	\end{figure}
}
\newcommand{\underlinestr}[2][10pt]{%
  \begingroup
  \fontsize{#1}{#1}\selectfont
  \ttfamily
  \def\temp{}%
  \count255=0
  \loop\ifnum\count255<#2
    \edef\temp{\temp\_}
    \advance\count255 1
  \repeat
  \temp
  \endgroup
}
\newcommand{\headerline}[3]{%
	\par\medskip\noindent
	\makebox[0pt][l]{#1}%
	\makebox[\textwidth][c]{#2}%
	\makebox[0pt][r]{\texttt{#3}}\par\medskip}
\newcommand{\HRule}{\rule{\linewidth}{0.3mm}} % Defines a new command for horizontal lines, change thickness here
\setcounter{page}{4} % Начало нумерации страниц
\regtotcounter{figure}
\regtotcounter{table}
\regtotcounter{equation}
\newtotcounter{bibitems}
\pretocmd{\bibitem}{\stepcounter{bibitems}}{}{}
\usepackage{tocbasic}

\DeclareTOCStyleEntries[
raggedentrytext,
linefill=\TOCLineLeaderFill,
indent=0pt,
numwidth=0pt,
numsep=0.5em,
dynnumwidth
]{tocline}{chapter,section}
\DeclareTOCStyleEntries[
raggedentrytext,
linefill=\TOCLineLeaderFill,
indent=1em,
dynnumwidth,
numsep=1ex
]{tocline}{subsection}
\DeclareTOCStyleEntries[
raggedentrytext,
linefill=\TOCLineLeaderFill,
indent=2em,
dynnumwidth,
numsep=0.5ex
]{tocline}{subsubsection}

\titlespacing*{\subsection}{0pt}{\baselineskip}{\baselineskip}
\titlespacing*{\subsubsection}{0pt}{\baselineskip}{\baselineskip}
\titlespacing*{\section}{0pt}{\baselineskip}{\baselineskip}
\titlespacing*{\equation}{0pt}{\baselineskip}{\baselineskip}